\chapter*{Цели и задачи}
\addcontentsline{toc}{chapter}{Введение}
\label{ch:intro_radio}

\textbf{Цель практической работы:} освоить методы имитационного моделирования
радиоканала подвижной связи: смоделировать траекторию перемещения абонента
(Mobile Station, MS) относительно базовой станции (Base Station, BS),
рассчитать принимаемый сигнал в комплексной низкочастотной форме
(комплексную огибающую), определить доплеровский сдвиг частоты, а также
исследовать статистические свойства сигнала при многолучевом распространении
(релеевские замирания).

\textbf{Задачи работы:}
\begin{enumerate}
    \item Смоделировать прямолинейную траекторию перемещения абонента
          с заданной скоростью и шагом дискретизации по расстоянию.
    \item Вычислить комплексную огибающую принимаемого сигнала для прямого
          луча BS--MS, определить амплитуду, фазу и доплеровский спектр.
    \item Смоделировать сценарий с большим числом рассеивателей
          (многолучевое распространение) и вычислить суммарный принимаемый
          сигнал.
    \item Построить и проанализировать характеристики сигнала: модуль, фазу,
          спектр Доплера, гистограмму и функцию распределения модуля,
          автокорреляционную функцию.
    \item Проверить соответствие распределения модуля принимаемого сигнала
          закону Рэлея.
    \item Оформить отчёт по выполненной работе.
\end{enumerate}

Моделирование выполняется на языке Python~3 с использованием библиотек
\texttt{numpy}, \texttt{matplotlib} и \texttt{scipy}.

\endinput